%%%%%%%%%%  scan idx 51  |  folio 38  %%%%%%%%%%
d'augmentation nulle), tandis que $\lambda_{-1}(N)$ est de filtration $q$ (formule (1.22)), il en résulte que $\gamma^n(N,x)$ est de filtration $j$, puisque en vertu de (1.18) le gradué associé à $A$ est ici \emph{intègre}. Utilisant les formules précédentes, on obtient immédiatement la forme indiquée pour $(\gamma^{q+j}(N,x))^{(j)}$.
\begin{flushright}cqfd.\end{flushright}

\section*{CHAPITRE II}
\begin{center}
Classes de faisceaux algébriques cohérents\par
et classes de Chern
\end{center}

\subsection*{\S 1. La théorie de Chow}

Dans tout ce chapitre, on suppose fixé un \emph{corps de base} $k$ qu'on supposera algébriquement clos, pour simplifier. Une grande partie de ce qui va suivre, et peut-être tout, est cependant valable sans cette restriction. Les espaces algébriques et les applications régulières considérés seront définis sur $k$.

Un espace algébrique est dit \emph{quasi-projectif} s'il est isomorphe à une partie localement fermée d'un espace projectif.

Soit $X$ une variété non singulière quasi-projective connexe de dimension $n$. On désigne par $A(X)$ et l'on appelle \emph{anneau de Chow} de $X$ l'anneau gradué des classes de cycles sur $X$ à équivalence rationnelle près. ($A^i(X)$ étant formé des classes de cycles de dimension $n-i$). Le fait que $A(X)$
